\scanpage{396}\mspage{135}

qui soit une $\mathrm{Hot}$-équivalence (p.ex. à l'aide d'une rétraction de $\unsure{\mathscr L|\mathcal L}_{n'}$ sur $\mathscr L_n$, \deleted{\unsure{$)_{\!2}$}} compatible avec sources et buts).

Ceci attire fortement l'attention sur le fait que si $Y,C$ \deleted{\unsure{deux}} deux catégories, la catégorie $\underline{\mathrm{Hom}}(Y,C)$ \uline{n'est pas} un représentant de \deleted{la catégorie} $\underline{\mathrm{Hom}}_{\mathrm{Hot}}(\mathrm{hot}(Y),\mathrm{hot}(C))$, on a seulement une \deleted{\unsure{appli}} flèche canonique

\[
\mathrm{hot}(\underline{\mathrm{Hom}}(Y,C))\longrightarrow \underline{\mathrm{Hom}}_{\mathrm{Hot}}(\mathrm{hot}(Y),\mathrm{hot}(C))
\tag*{(*)}
\]

i.e.

\[
\underbrace{\unsure{\mathrm{hom}|\mathrm{hot}}(\underline{\mathrm{Hom}}(Y,C))\times\mathrm{hot}(Y)}_{\substack{\Vert\\ \mathrm{hot}(\underline{\mathrm{Hom}}(Y,C)\times Y)}}\longrightarrow\mathrm{hot}(C)
\]

laquelle flèche est simplement déduite de

\[
\underline{\mathrm{Hom}}(Y,C)\times Y\longrightarrow C.
\]

Mais la flèche (*) n'a aucune envie d'être iso. C'est parce qu'elle ne l'est \inserted{déjà} pas dans le cas particulier $Y=\Delta^1$, essentiellement, qui nous oblige, pour faire une théorie des chemins \inserted{dans \deleted{\unsure{le}}\unsure{les} catégories} qui soit topologiquement raisonnable, de ne pas nous borner à $\Delta^1$ comme types de chemins, mais à prendre des intervalles \uline{composés} $\tau\unsure{\in|+}\unsure{\mathrm{Ty}|\underline{\mathrm{Ty}}}$.

% archive sheet number: 395
